% GENERATED FILE. Edit canonical lessons, then run npm run generate:books.
% canonical-source-hash: fnv1a64:3aed480a5710f33f
% canonical-lessons: RU-C445-fonar, RU-C445-skameyka, RU-C445-urna, RU-C445-kiosk, RU-C445-klumba

\chapter{Lamps and benches}
\label{ch:ru-lamps-and-benches}
\hlchaptermodality{445}

\begin{chapteropening}
\textbf{By the end of this chapter:} \emph{I can describe the street: street lamps, benches, litter bins, kiosks and flower beds.}

Complete the last lesson of chapter 445: I can describe the street: street lamps, benches, litter bins, kiosks and flower beds.
\end{chapteropening}

\section[\texorpdfstring{\ru{фонарь} (fonár)}{фонарь (fonár)}]{\ru{фонарь} (fonár) — a street lamp; a torch, a lantern}

\label{lesson:RU-C445-fonar}



\begin{quote}
Before the new one: say the Russian for a fortress, then the Russian for a city hall.
\end{quote}

\subsection*{You'll want to know: \ru{фонарь}}
\textbf{\ru{фонарь}} (\emph{fonár}) — \textquotedblleft{}a street lamp; a torch, a lantern\textquotedblright{}.

\begin{quote}
\ru{Вечером} \ru{на} \ru{улице} \ru{зажигаются} \ru{фонари}.
\end{quote}

\begin{quote}
\emph{Vécherom na úlitse zazhigáyutsya fonarí.}
\end{quote}

\begin{quote}
In the evening the street lamps come on.
\end{quote}

Masculine; the stress moves: \emph{\ru{фонаря}} (\emph{fonaryá}), \emph{\ru{фонари}} (\emph{fonarí}). \emph{\ru{Карманный} \ru{фонарь}} — a torch.

\begin{cousinweb}
\emph{\ru{фонарь}} goes back to Byzantine Greek \emph{phanárion} 'a lamp, a torch' (EtymDB).
\end{cousinweb}

\begin{grammarlens}[title={What you've built}]
One more word for telling, explaining and describing at B1.
\end{grammarlens}

\subsection*{Your turn}
\begin{itemize}
\raggedright
  \item \emph{Say it:} \emph{fonár}
  \item \emph{Say it:} \emph{fonár}, once more
  \item \emph{Say it:} a sentence of your own with \emph{fonár}
  \item \emph{Recall:} say the Russian for a fortress, then the Russian for a city hall, then say \emph{fonár} again
\end{itemize}

\begin{tcolorbox}[breakable,colback=teal!4,colframe=teal!35!black,title={Before you move on}]
What does \ru{фонарь} mean? (\textquotedblleft{}A street lamp; a torch, a lantern\textquotedblright{}.) Say it once more.
\end{tcolorbox}
\section[\texorpdfstring{\ru{скамейка} (skaméyka)}{скамейка (skaméyka)}]{\ru{скамейка} (skaméyka) — a bench}

\label{lesson:RU-C445-skameyka}



\begin{quote}
Before the new one: say the Russian for a city hall, then the Russian for a street lamp.
\end{quote}

\subsection*{You'll want to know: \ru{скамейка}}
\textbf{\ru{скамейка}} (\emph{skaméyka}) — \textquotedblleft{}a bench\textquotedblright{}.

\begin{quote}
\ru{Бабушки} \ru{сидят} \ru{на} \ru{скамейке} \ru{у} \ru{подъезда}.
\end{quote}

\begin{quote}
\emph{Bábushki sidyát na skaméyke u podyézda.}
\end{quote}

\begin{quote}
The grandmothers are sitting on the bench by the entrance.
\end{quote}

Feminine; a diminutive of \emph{\ru{скамья}}.

\begin{cousinweb}
\emph{\ru{скамейка}} is a diminutive of \emph{\ru{скамья}} 'a bench', which EtymDB takes from Byzantine Greek \emph{skamnía}. EtymDB lists \emph{\ru{скамейка}} itself without a source.
\end{cousinweb}

\begin{grammarlens}[title={What you've built}]
One more word for telling, explaining and describing at B1.
\end{grammarlens}

\subsection*{Your turn}
\begin{itemize}
\raggedright
  \item \emph{Say it:} \emph{skaméyka}
  \item \emph{Say it:} \emph{skaméyka}, once more
  \item \emph{Say it:} a sentence of your own with \emph{skaméyka}
  \item \emph{Recall:} say the Russian for a city hall, then the Russian for a street lamp, then say \emph{skaméyka} again
\end{itemize}

\begin{tcolorbox}[breakable,colback=teal!4,colframe=teal!35!black,title={Before you move on}]
What does \ru{скамейка} mean? (\textquotedblleft{}A bench\textquotedblright{}.) Say it once more.
\end{tcolorbox}
\section[\texorpdfstring{\ru{урна} (úrna)}{урна (úrna)}]{\ru{урна} (úrna) — a litter bin; an urn}

\label{lesson:RU-C445-urna}



\begin{quote}
Before the new one: say the Russian for a street lamp, then the Russian for a bench.
\end{quote}

\subsection*{You'll want to know: \ru{урна}}
\textbf{\ru{урна}} (\emph{úrna}) — \textquotedblleft{}a litter bin; an urn\textquotedblright{}.

\begin{quote}
\ru{Он} \ru{бросил} \ru{старую} \ru{газету} \ru{в} \ru{урну}.
\end{quote}

\begin{quote}
\emph{On brósil stáruyu gazétu v úrnu.}
\end{quote}

\begin{quote}
He threw the old newspaper into the bin.
\end{quote}

Feminine. \emph{\ru{Урна} \ru{для} \ru{голосования}} — a ballot box.

\begin{cousinweb}
\emph{\ru{урна}} was borrowed from Latin \emph{urna} 'a vessel, a jar' (EtymDB).
\end{cousinweb}

\begin{grammarlens}[title={What you've built}]
One more word for telling, explaining and describing at B1.
\end{grammarlens}

\subsection*{Your turn}
\begin{itemize}
\raggedright
  \item \emph{Say it:} \emph{úrna}
  \item \emph{Say it:} \emph{úrna}, once more
  \item \emph{Say it:} a sentence of your own with \emph{úrna}
  \item \emph{Recall:} say the Russian for a street lamp, then the Russian for a bench, then say \emph{úrna} again
\end{itemize}

\begin{tcolorbox}[breakable,colback=teal!4,colframe=teal!35!black,title={Before you move on}]
What does \ru{урна} mean? (\textquotedblleft{}A litter bin; an urn\textquotedblright{}.) Say it once more.
\end{tcolorbox}
\section[\texorpdfstring{\ru{киоск} (kiósk)}{киоск (kiósk)}]{\ru{киоск} (kiósk) — a kiosk, a news stand}

\label{lesson:RU-C445-kiosk}



\begin{quote}
Before the new one: say the Russian for a bench, then the Russian for a litter bin.
\end{quote}

\subsection*{You'll want to know: \ru{киоск}}
\textbf{\ru{киоск}} (\emph{kiósk}) — \textquotedblleft{}a kiosk, a news stand\textquotedblright{}.

\begin{quote}
\ru{В} \ru{киоске} \ru{продают} \ru{газеты} \ru{и} \ru{билеты}.
\end{quote}

\begin{quote}
\emph{V kióske prodayút gazéty i biléty.}
\end{quote}

\begin{quote}
The kiosk sells newspapers and tickets.
\end{quote}

Masculine.

\begin{cousinweb}
\emph{\ru{киоск}} came from French \emph{kiosque} (EtymDB), which goes back through Turkish \emph{köşk} 'a pavilion' to Persian.
\end{cousinweb}

\begin{grammarlens}[title={What you've built}]
One more word for telling, explaining and describing at B1.
\end{grammarlens}

\subsection*{Your turn}
\begin{itemize}
\raggedright
  \item \emph{Say it:} \emph{kiósk}
  \item \emph{Say it:} \emph{kiósk}, once more
  \item \emph{Say it:} a sentence of your own with \emph{kiósk}
  \item \emph{Recall:} say the Russian for a bench, then the Russian for a litter bin, then say \emph{kiósk} again
\end{itemize}

\begin{tcolorbox}[breakable,colback=teal!4,colframe=teal!35!black,title={Before you move on}]
What does \ru{киоск} mean? (\textquotedblleft{}A kiosk, a news stand\textquotedblright{}.) Say it once more.
\end{tcolorbox}
\section[\texorpdfstring{\ru{клумба} (klúmba)}{клумба (klúmba)}]{\ru{клумба} (klúmba) — a flower bed}

\label{lesson:RU-C445-klumba}



\begin{quote}
Before the new one: say the Russian for a litter bin, then the Russian for a kiosk.
\end{quote}

\subsection*{You'll want to know: \ru{клумба}}
\textbf{\ru{клумба}} (\emph{klúmba}) — \textquotedblleft{}a flower bed\textquotedblright{}.

\begin{quote}
\ru{Перед} \ru{школой} \ru{большая} \ru{клумба} \ru{с} \ru{розами}.
\end{quote}

\begin{quote}
\emph{Péred shkóloy bolsháya klúmba s rózami.}
\end{quote}

\begin{quote}
There is a big flower bed with roses in front of the school.
\end{quote}

Feminine.

\begin{cousinweb}
\emph{\ru{клумба}} is usually traced to English \emph{clump} 'a cluster of plants'. EtymDB lists the word beside a Polish cousin, \emph{klomb}, but gives no source.
\end{cousinweb}

\begin{grammarlens}[title={What you've built}]
One more word for telling, explaining and describing at B1.
\end{grammarlens}

\subsection*{Your turn}
\begin{itemize}
\raggedright
  \item \emph{Say it:} \emph{klúmba}
  \item \emph{Say it:} \emph{klúmba}, once more
  \item \emph{Say it:} a sentence of your own with \emph{klúmba}
  \item \emph{Recall:} say the Russian for a litter bin, then the Russian for a kiosk, then say \emph{klúmba} again
\end{itemize}

\begin{tcolorbox}[breakable,colback=teal!4,colframe=teal!35!black,title={Before you move on}]
What does \ru{клумба} mean? (\textquotedblleft{}A flower bed\textquotedblright{}.) Say it once more.
\end{tcolorbox}
